\PassOptionsToPackage{usenames,dvipsnames}{xcolor}
\documentclass[12pt, a4paper]{article}
\usepackage{tabularx}
\usepackage[export]{adjustbox}
\usepackage[brazil]{babel}
\usepackage{lipsum}

\usepackage[utf8x]{inputenc}
\usepackage[lmargin=2cm,tmargin=2cm,rmargin=2cm,bmargin=2cm]{geometry}

% \makeatletter
% \patchcmd{\pgf@circ@zigzag}{\pgfsetbeveljoin}{\pgfsetroundjoin}
%     {\typeout{Switching to safe resistors!}}
%     {\typeout{Patching resistors failed}}
% \makeatother

\usepackage[theorems]{tcolorbox}
% \usepackage[most]{tcolorbox}

\usepackage{cancel}
\usepackage{amsmath}


\definecolor{disp1}{HTML}{FFD723}
\definecolor{disp2}{HTML}{FF00DD}
\definecolor{disp3}{HTML}{FFD723}
\usepackage[american]{circuitikz}

\usepackage{lmodern,cmap}
\usepackage[bottom]{footmisc}
\usepackage{tikz, tkz-euclide}
\usepackage{relsize}
\usetikzlibrary{arrows, positioning}

\usepackage{graphicx}
\usepackage{mathtools}
\usepackage{amssymb}
\definecolor{line}{HTML}{005f7f}
\definecolor{disp}{HTML}{329881}
\usetikzlibrary{bending}

\makeatletter
\ctikzset{current arrow color/.initial=black}% create key

\let\old@circ@drawcurrent=\pgf@circ@drawcurrent
\def\pgf@circ@drawcurrent{\old@circ@drawcurrent}

\pgfdeclareshape{currarrow}{
 \anchor{center}{
  \pgfpointorigin
 }
 \anchor{tip}{
  \pgfpointorigin
  \pgf@circ@res@step = \pgf@circ@Rlen
  \divide \pgf@circ@res@step by 16
  \pgf@x  =\pgf@circ@res@step
 }
 \behindforegroundpath{

  \pgfscope
  \pgf@circ@res@step = \pgf@circ@Rlen
  \divide \pgf@circ@res@step by 16

  \pgfpathmoveto{\pgfpoint{-.7\pgf@circ@res@step}{0pt}}
  \pgfpathlineto{\pgfpoint{-.7\pgf@circ@res@step}{-.8\pgf@circ@res@step}}
  \pgfpathlineto{\pgfpoint{1\pgf@circ@res@step}{0pt}}
  \pgfpathlineto{\pgfpoint{-.7\pgf@circ@res@step}{.8\pgf@circ@res@step}}
  \pgfpathlineto{\pgfpoint{-.7\pgf@circ@res@step}{0pt}}
  \pgfsetcolor{\pgfkeysvalueof{/tikz/circuitikz/current arrow color}}
  \pgfusepath{draw,fill}

  \endpgfscope

 }
}
\pgfdeclareshape{flowarrow}{
 \anchor{center}{\pgfpointorigin}
 \anchor{tip}{
  \pgfpointorigin
  \pgf@circ@res@step = \pgf@circ@Rlen
  \divide \pgf@circ@res@step by 16
  \pgf@x  =\pgf@circ@res@step
 }
 \behindforegroundpath{
  \pgfscope
  \pgf@circ@res@step = \pgf@circ@Rlen
  \divide \pgf@circ@res@step by 4
  \pgfpathmoveto{\pgfpoint{-\pgf@circ@res@step}{0pt}}
  \pgfpathlineto{\pgfpoint{\pgf@circ@res@step}{0pt}}
  %\pgfsetcolor{\pgfkeysvalueof{/tikz/circuitikz/color}}
  \pgfsetcolor{\pgfkeysvalueof{/tikz/circuitikz/current arrow color}}
  \pgfusepath{draw}
  \pgftransfoRMShift{\pgfpoint{\pgf@circ@res@step}{0pt}}
  \pgf{currarrow}{tip}{}{}{\pgfusepath{fill}}
  \endpgfscope
 }
}
\makeatother

\usepackage{multirow}

\usepackage{array}

\usepackage{indentfirst}
\usepackage{xurl}
\usepackage{graphicx}
\usepackage{ragged2e}

\usepackage[font=footnotesize,labelfont=bf]{caption}

% \usepackage{fontspec}

\usepackage[hidelinks]{hyperref}
\usepackage{nameref}

\captionsetup{belowskip=3.241pt,aboveskip=3.241pt, position=top}

\usepackage[titles]{tocloft}
\usepackage{titlesec}

\usepackage{microtype}

\makeatletter
\renewcommand{\@tocrmarg}{5em}
\makeatother

\usepackage{titletoc}
% tkz-euclide ja carregado acima (linha 31); nao recarregar para evitar conflito com pgfplots
\captionsetup{belowskip=0pt,aboveskip=0pt}
\usepackage{tikz-3dplot}
\usetikzlibrary{arrows}
\usepackage{siunitx}
\usepackage{mathtools}
\usepackage{amssymb}
\usepackage{newfloat}

% pgfplots carregado aqui para compatibilidade com tkz-euclide
\usepackage{pgfplots}
\pgfplotsset{compat=1.18}
\usepgfplotslibrary{statistics, fillbetween}

\pgfplotsset{
	expaxis/.style={
		grid=major,
		grid style={gray!25},
		tick label style={font=\footnotesize},
		label style={font=\small},
		legend cell align=left,
		legend style={
			font=\footnotesize,
			draw=black!35,
			fill=white,
			fill opacity=0.92,
			text opacity=1,
			rounded corners=1.5pt,
			inner xsep=5pt,
			inner ysep=4pt,
			row sep=1pt,
		},
	},
	explegendNE/.style={
		legend style={at={(axis description cs:0.975,0.975)},anchor=north east},
	},
	explegendNW/.style={
		legend style={at={(axis description cs:0.025,0.975)},anchor=north west},
	},
	explegendSE/.style={
		legend style={at={(axis description cs:0.975,0.025)},anchor=south east},
	},
	explegendSW/.style={
		legend style={at={(axis description cs:0.025,0.025)},anchor=south west},
	},
}

\usetikzlibrary{positioning}

\usepackage{sectsty}


\DeclareFloatingEnvironment[fileext=grf,placement={!htbp},
 name=Circuito,listname={CIRCUITOS},]{circuito}


\DeclareFloatingEnvironment[fileext=grf,placement={!htbp},name=Gráfico,listname={GRÁFICOS},]{grafico}

% \sectionfont{\fontsize{12pt}{1.5}\selectfont}
%
% \subsectionfont{\fontsize{12pt}{1.5}\selectfont}
%
% \subsubsectionfont{\fontsize{12pt}{1.5}\selectfont}

\newcommand{\tikzAngleOfLine}{\tikz@AngleOfLine}
\def\tikz@AngleOfLine(#1)(#2)#3{%
 \pgfmathanglebetweenpoints{%
  \pgfpointanchor{#1}{center}}{%
  \pgfpointanchor{#2}{center}}
 \pgfmathsetmacro{#3}{\pgfmathresult}%
}

\makeatletter
\newenvironment{system}
{%
\let\@ifnextchar\new@ifnextchar
\def\arraystretch{1.2}%
\left\lbrace
\array{@{}l@{}}%
}
{%
\endarray
\right.\kern-\nulldelimiterspace
}
\makeatother

\usepackage{enumitem}

\setlist{nosep}

\usepackage{import}
\usepackage{xifthen}
\usepackage{pdfpages}
\usepackage{transparent}
\usepackage{booktabs}
\usepackage[onehalfspacing]{setspace}

\newenvironment{citacaolonga}{\vspace{1em}\par\hfill\noindent\begin{minipage}[c]{8cm}\setstretch{1.0}\footnotesize}{\end{minipage}\vspace{.75em}}

\setlength{\parindent}{1.25cm}

\makeatletter
\renewcommand\section{%
  \@startsection{section}{1}
                {\z@}%
                {1em}%
                {1em}
                {\Large\bfseries}%
}

\makeatletter
\renewcommand\subsection{%
  \@startsection{subsection}{2}{\z@}{1em}%
  {1em}{\normalfont\large\bfseries}
  }
\makeatother

\makeatletter
\renewcommand\subsubsection{%
  \@startsection{subsubsection}{3}{\z@}{1em}%
  {1em}{\normalfont}
  }
\makeatother

\usepackage{listings}

\definecolor{background}{RGB}{46, 52, 64}
\definecolor{string}{RGB}{255, 127, 0}
\definecolor{comment}{RGB}{95, 95, 128}
\definecolor{normal}{RGB}{248, 248, 242}
\definecolor{identifier}{RGB}{166, 226, 46}
\definecolor{bblue}{RGB}{0, 141, 254}

% \renewcommand{\rmdefault}{cmbr}

\lstset{
 % language=c,                			% choose the language of the code
 % numbers=left,                   		% where to put the line-numbers
 stepnumber=1,                   		% the step between two line-numbers.
 numbersep=5pt,                  		% how far the line-numbers are from the code
 numberstyle=\color{fundd}\footnotesize\bfseries\ttfamily, 		% choose the background color. You must add \usepackage{color}
 showspaces=false,               		% show spaces adding particular underscores
 showstringspaces=false,         		% underline spaces within strings
 showtabs=false,                 		% show tabs within strings adding particular underscores
 tabsize=2,             		% sets default tabsize to 2 spaces
 % backgroundcolor=\color{black!10},
 captionpos=t,                   		% sets the caption-position to bottom
 breaklines=true,                		% sets automatic line breaking
 breakatwhitespace=false,         		% sets if automatic breaks should only happen at whitespace
 inputencoding=utf8,
 escapechar={|},
 % title=\lstname,                 		% show the filename of files included with \lstinputlisting;
 basicstyle=\color{fundd}\ttfamily,					% sets font style for the code
 keywordstyle=\color{bblue}\bfseries\ttfamily,	% sets color for keywords
 stringstyle=\color{magenta}\bfseries\ttfamily,		% sets color for strings
 commentstyle=\color{comment!75}\ttfamily,	% sets color for comments
 emph={},
 % frame=tlbr,
 % frame=lr,
 % framesep=8pt,
 % framerule=0pt,
 emphstyle=\color{identifier}\ttfamily,
language=c,   % R code
literate=
{á}{{\'a}}1
{à}{{\`a}}1
{ã}{{\~a}}1
{é}{{\'e}}1
{ê}{{\^e}}1
{í}{{\'i}}1
{ó}{{\'o}}1
{õ}{{\~o}}1
{ú}{{\'u}}1
{ü}{{\"u}}1
{ç}{{\c{c}}}1
}

\renewcommand\bfdefault{b}
